\section{Enunciado bajo las entradas de la vía 2}
Sean $K=\lceil\sqrt n\rceil$, $j=\lceil tK\rceil\le M=\lceil\delta n\rceil$, $s=1/j$.
La esperanza subprobabilística de llegada es $\E_*Z=\E_{\bar\nu}[Z;G_j,S_M]$.
Todos los índices de línea y las coordenadas de las casillas se miden en el tablero
original. Al retener filas y columnas no se vuelven a definir las diagonales.
En una máscara basal $A$, sea $e(Y)=\log(Q_Y/\per A)-(L_Y-L_T)$: se activan las
capacidades diagonales sólo entre las marcas de $Y$, manteniendo disponibles para
las torres todos los sitios de $A$. Para $S$ de $v$ sitios distintos, ponemos
\[
 [S]e=\sum_{Y\subseteq S}(-1)^{v-|Y|}e(Y).
\]
Llamamos $V$ al error diagonal tras restar las líneas individuales y los incrementos
de dos conflictos en líneas distintas, con las mismas convenciones que en la entrega previa.
Una sustracción de dos conflictos sólo tiene proyección sobre la unión de sus sitios.
Una sustracción de línea individual sólo puede sobrevivir si todo $S$ está en esa línea.

\paragraph{Hipótesis uniformes admitidas.}
Se usan la difusión y la brecha de torres sobre los tableros considerados, la estabilidad
de los menores al retirar un número fijo de filas y columnas y la cota de prescripción
de llegada $C_h(j/n)^h$ para $h$ coordenadas originales distintas. Se usa la cota agregada
por exceso de R1 y R2, con las colas en grado y en la mezcla gamma que se precisan en la
parte 2. Las identidades de valor, Schur y caras son las de la entrega sobre grupos
interlíneas, partes 1--4, sobre el mismo exterior $D_0=A\setminus S$. El radio complejo
de las actividades locales puede depender de $v$, pero es independiente de $j$.
No se exige que las activaciones de tamaño creciente estén en ese radio.

\paragraph{Resultado.}
Para cada $v$ fijado antes de tomar límites, en toda la banda y para $j$ suficientemente
grande en función de $v$, se cumple
\[
 \E_*\sum_{\substack{S\subseteq A\\ |S|=v}}|[S]V|\ \le\ C_v/n .
\]
La prueba incluye sitios que comparten filas o columnas y componentes con ciclos.
En particular, todo corte finito de la suma es $o((\log t)/n)$ uniformemente al crecer $t$.
El resultado amplía el de tres componentes forestales con coordenadas RC distintas;
no supone ninguna estimación del crecimiento de $C_v$.

\paragraph{Qué grafo determina la cuenta.}
Se usan las componentes del grafo bipartito que une cada sitio con sus líneas
de fila, columna y ambas diagonales compartidas dentro de $S$. Las líneas con un
solo sitio se pueden dejar como hojas. Denotamos por $c$ el número de componentes
de ese grafo, incluidas las que en una activación parcial puedan dividirse.
Es una partición fija, común a todos los $Y\subseteq S$. Los sitios sin ningún
vecino diagonal se eliminan de la discusión: su activación no cambia las capacidades
y su proyección es cero. En particular toda contribución no nula tiene una igualdad
diagonal entre sitios distintos.
